\section{The mailbox relay}
\label{sec:mailbox-relay}

The hosted mailbox carries sealed \file{.kqpb} envelopes and public-tree
slices. It indexes packages by the recipient fingerprint in the outer
header and never unseals them; wrapped shares and private keys stay on the
device. Your provider gives you a URL and an API key --- you do not run the
mailbox or mint keys yourself. Official \code{loadkey} / \code{relay}
commands authenticate that host with a KeyQuorum-signed provider
certificate before sending a bearer.

\begin{lstlisting}
export KEYQUORUM_RELAY_URL=https://relay.example.com
# Load once; omit the key so it is prompted (stays out of shell history).
keyquorum loadkey --url https://relay.example.com
keyquorum relay push --dir ./bridge-packages
keyquorum relay pull --output-dir ./inbox
keyquorum --db alice.sqlite relay pull --import --share-file alice.key
# --api-key still works for scripts and is stored after a successful check.
keyquorum relay push --dir ./bridge-packages --api-key "$PUSH_KEY"
\end{lstlisting}

\code{loadkey} first verifies the relay's KeyQuorum-signed identity, then
checks the key with the mailbox, then stores the key hash and a sealed copy
of the bearer in the personal database. Later \code{relay} / \code{tree
fetch} commands repeat that identity check and re-check the hash before
using the key. An optional signed revocation list can be pointed at with
\code{KEYQUORUM_PROVIDER_KRL}.

\code{relay push} also uploads every public tree in \flag{--db} and
\textbf{merges} it into the mailbox (unrelated nodes stay put). \code{relay
pull} merges the returned slices into \flag{--db} (then \flag{--import}
opens envelopes). \code{tree fetch} syncs topology without an envelope.
Remote mailboxes must be \code{https://}; \code{http://} is accepted only
for loopback.

\subsection{The device mailbox}

Device copy, move, and relocate (Section~\ref{sec:device-transfer}) use a
second mailbox on the same host. \code{device.push} stores a sealed letter
and publishes a device's public descriptor (device id, verify key, and slot
public keys, signed by the device key). \code{device.pull} is bound to one
recipient fingerprint, the same way \code{inbox.pull} is, and reads only
that fingerprint's letters. A missing bearer is rejected, and a key of the
wrong scope is rejected --- the host does not mint keys over HTTP. Your
provider issues \code{device.push} and \code{device.pull} the same way it
issues inbox keys; \code{loadkey} stores them.

The host refuses a raw \code{KQTX} package and never unseals a letter.
Bridge import does not read this mailbox, and device letters are not
accepted into the bridge inbox.

\subsection{Recovering from a lost key or a missed pull}

If a key is lost or rotated, load the replacement your provider issues.
Envelopes already delivered are unchanged.

\begin{lstlisting}
keyquorum relay pull --after <id> --output-dir ./inbox
# or, to unseal them into --db directly:
# keyquorum relay pull --after <id> --import --share-file alice.key
\end{lstlisting}

replays anything not yet downloaded; \code{bridge private import} still
rejects stale generations, so replaying an old envelope by mistake changes
nothing.
